%% sections/supp.tex -- OUTLINE and first draft of the supplementary material.
%% The supplement is a separate, anonymized PDF (due Mon Nov 23, 2026 AoE); main.tex does NOT input this file.
%% Intended use: a supp.tex driver with the same preamble/macros that inputs this file (references to the main
%% paper are written as "main paper, Sec. X" because labels of the main PDF are not available here).
%% Content marked \TODO is planned but not written yet; every number must come from a logged run.

\section{Derivation of the gap smoother}
\label{supp:smoother}

\paragraph{Setting.}
Consider one Kalman filter of the predictor (one residual element, or one DCT band whose coefficients share the
covariance). Let $a<b$ be consecutive computed steps, $\mathcal Y_a$ the observations up to step $a$, and $a<i<b$ a
skipped step. No observation falls into the gap, so the predicted moments are propagated without update:
$\prior{\kstate}_{i}=\kA_i\prior{\kstate}_{i-1}$ and $\prior{\kP}_i=\kA_i\prior{\kP}_{i-1}\kA_i^\top+\kQ_i$ for
$a+1<i\le b$, starting from $\prior{\kstate}_{a+1}=\kA_{a+1}\post{\kstate}_a$. Define the transition
$\kPhi(b,i)=\kA_b\kA_{b-1}\cdots\kA_{i+1}$, with $\kPhi(b,b)=\mathbf I$, and the innovation variance
$S_b=\mathbf h^\top\prior{\kP}_b\mathbf h+R$.

\paragraph{Proposition.}
For $a<i\le b$,
\begin{equation}
  \E\!\left[\kstate_i\mid\mathcal Y_a,z_b\right]=\prior{\kstate}_i+\prior{\kP}_i\kPhi(b,i)^\top\mathbf h\,S_b^{-1}e_b,
  \qquad e_b=z_b-\mathbf h^\top\prior{\kstate}_b,
  \label{eq:supp-smooth}
\end{equation}
and the first component gives the gain $g_i=[\prior{\kP}_i\kPhi(b,i)^\top]_{00}/S_b$ of the main paper.

\paragraph{Proof by conditioning.}
Iterating the state equation, $\kstate_b=\kPhi(b,i)\kstate_i+\mathbf w_{i\to b}$, where $\mathbf w_{i\to b}$ is a sum of
process-noise terms of steps $i+1,\dots,b$; it has zero mean and is independent of $\kstate_i$ and of $\mathcal Y_a$.
With $z_b=\mathbf h^\top\kstate_b+\eta_b$ and $\eta_b\sim\mathcal N(0,R)$ independent of everything else,
$\mathrm{Cov}(\kstate_i,z_b\mid\mathcal Y_a)=\prior{\kP}_i\kPhi(b,i)^\top\mathbf h$ and
$\mathrm{Var}(z_b\mid\mathcal Y_a)=S_b$. Conditioning the jointly Gaussian pair $(\kstate_i,z_b)$ given $\mathcal Y_a$
on $z_b$ gives \eqref{eq:supp-smooth}.

\paragraph{Equivalence with the RTS smoother.}
The RTS recursion~\cite{rauch1965maximum,sarkka2013bayesian} over $[a,b]$, using the observations up to $b$, reads
$\smooth{\kstate}_i=\post{\kstate}_i+\mathbf C_i(\smooth{\kstate}_{i+1}-\prior{\kstate}_{i+1})$ with
$\mathbf C_i=\post{\kP}_i\kA_{i+1}^\top(\prior{\kP}_{i+1})^{-1}$, initialized with the filtered estimate
$\smooth{\kstate}_b=\post{\kstate}_b=\prior{\kstate}_b+\prior{\kP}_b\mathbf h S_b^{-1}e_b$. Inside the gap,
$\post{\kstate}_i=\prior{\kstate}_i$ and $\post{\kP}_i=\prior{\kP}_i$. We show by backward induction that
$\smooth{\kstate}_i-\prior{\kstate}_i=\prior{\kP}_i\kPhi(b,i)^\top\mathbf h S_b^{-1}e_b$. The case $i=b$ holds because
$\kPhi(b,b)=\mathbf I$. If it holds for $i+1$, then
\begin{equation*}
  \smooth{\kstate}_i-\prior{\kstate}_i
  =\prior{\kP}_i\kA_{i+1}^\top(\prior{\kP}_{i+1})^{-1}\prior{\kP}_{i+1}\kPhi(b,i+1)^\top\mathbf h S_b^{-1}e_b
  =\prior{\kP}_i\big(\kPhi(b,i+1)\kA_{i+1}\big)^\top\mathbf h S_b^{-1}e_b,
\end{equation*}
and $\kPhi(b,i+1)\kA_{i+1}=\kPhi(b,i)$. \hfill$\square$

\paragraph{Consequences.}
\begin{itemize}
  \item At the newest observation the smoothed and filtered estimates coincide, so \pcgs leaves the predictor for the
  next gap unchanged (main paper, Remarks of the \pcgs section).
  \item \emph{Random walk.} For the model without trend ($\kA=1$, $\kQ=q\ds$), an exact observation at $a$
  ($\post{\kP}_a=0$) and $R\to0$, the gain is $g_i=(s_i-s_a)/(s_b-s_a)$, \ie, linear interpolation in the time
  coordinate. With the reuse forecast $\hat\res_i=\res_a$ this is the linear interpolation between $\res_a$ and
  $\res_b$ used for \seacache + \pcgs in the ablations.
  \item \emph{DCT variant.} Each band $\ell$ has its own $\kA$, $\kQ$, $\kP$ and hence its own gains $g_{i,\ell}$. Band
  normalization cancels in \eqref{eq:supp-smooth}, so the correction in residual space is
  $\resf_i=\hat\res_i+\mathcal D^{-1}\big(\mathbf G_i\odot\mathcal D(\res_b-\hat\res_b)\big)$ with the per-channel
  orthonormal 2D DCT $\mathcal D$ and the gain map $\mathbf G_i$ (value $g_{i,\ell}$ on band $\ell$).
  \item \emph{Gain shape.} \TODO{Figure: $g_i$ per band for the modal 13-NFE schedule of \flux
  (0,1,2,4,6,9,13,18,25,33,41,47,49). Preliminary computation with the \flux calibration and $\kappa=1$ (no online
  adaptation): within the seven-step gaps the gain rises monotonically from 0.04--0.19 at the first skipped step to
  0.74--0.88 at the last one (high-frequency bands more convex); recompute with the logged $\kappa$.}
  \item \emph{First-order optimality.} \TODO{Write out the linearized error propagation
  $\lat_b-\lat^{\mathrm{ref}}_b\approx\sum_i\mathbf M_{b,i}(\sigma_{i+1}-\sigma_i)\mathbf J_i(\hat\res_i-\res_i)$ and the
  conditional-mean argument of the main paper.}
\end{itemize}

\section{Forecaster details}
\label{supp:forecaster}
\begin{itemize}
  \item DCT bands: orthonormal DCT-II per channel on the $64{\times}64$ token grid; seven radial bands with normalized
  radius $r=\sqrt{(u/H)^2+(v/W)^2}/\sqrt2$ and edges $\{0,\tfrac1{64},\tfrac1{32},\tfrac1{16},\tfrac18,\tfrac14,
  \tfrac12,1\}$; band normalization by the calibrated residual power.
  \item Time coordinate (following SpecKCast~\cite{rajabi2026speckcast}): $\lambda_i=2\log((1-\sigma_i)/\sigma_i)$
  (clamped), $t_\ell(i)=\mathrm{sigmoid}((\lambda_i+\log p_\ell)/T_w)$ with $T_w=2$, rescaled to $[-1,1]$; $p_\ell$ is
  the calibrated power of the clean-sample estimate in band $\ell$. Per-element variant: one band with the
  count-weighted mean power.
  \item Calibration (16 COCO captions, disjoint from the 64 development and the 200 DrawBench prompts): residual band
  powers, $p_\ell$ from the last 25 steps, process-noise intensity $q_\ell$ as the median over steps of the mean squared
  change of the finite-difference trend per unit warped time, initial trend variance $2q_\ell$.
  \item Initialization at the first computed step: $m=z_0$, $\nu=0$, $\kP=\mathrm{diag}(R,2q)$.
  \item Process-noise adaptation: $\gamma_b=\mathrm{clip}_{[0.1,10]}(\overline{e_b^2}/S_b)$ with the band-mean squared
  innovation, $\kappa\gets\mathrm{clip}_{[0.1,10]}(\kappa\gamma_b^{0.2})$.
  \item Calibration-free variant: unit band powers, $q=7400$, initial trend variance $15000$, uniform step time.
  \TODO{state which variant is used for video.}
\end{itemize}

\section{Implementation of \pcgs}
\label{supp:impl}
\begin{itemize}
  \item Bookkeeping per skipped step: latent $\lat_i$ (only $\lat_{a+1}$ is needed for the exact replay), conditioning
  $\cond_i$, forecast $\hat\res_i$, and the prior covariance of every band; memory cost (\flux): \TODO{MiB}.
  \item Exact replay: same tensors and arithmetic as the diffusers flow-matching Euler step (fp32 accumulation of the
  bf16 product, cast to bf16); with the identity smoother ($\resf_i=\hat\res_i$) the replay reproduces the uncorrected
  run bitwise (unit test).
  % provenance: research_v2/ablations/ABLATIONS_tests.md, test U3; T1 for the next item.
  \item The implemented gains match brute-force joint-Gaussian conditioning of the damped-trend model to a relative
  error below $10^{-15}$ for both parametrizations (unit test T1, float64, synthetic residuals, \flux calibration).
  \item Continuation at $b$ and the gate refresh (first-block adaptive normalization and \sea filter); the last step is
  corrected but needs no refresh.
  \item NFE accounting: forward pre-hooks on all 19+38 blocks, cross-checked against the forward pass (zero mismatches,
  zero partial block stacks in every run).
  \item HunyuanVideo port: Euler sampler, single guidance-distilled branch; the official gate binds the scheduler that
  exists at load time, so its \sea filter reads $\sigma=1-i/1000$ instead of the sampling schedule. We reproduce this
  behavior for all gated methods (\seacache's thresholds were tuned with it) and use the live schedule for the
  re-integration. \TODO{confirm with the final port.}
  \item Wan2.1 port: UniPC multistep solver with classifier-free guidance; per-branch filters; solver-state snapshot at
  the start of every gap and replay of the gap with the corrected outputs. \TODO{validate; report the bitwise test.}
\end{itemize}

\section{Evaluation protocol}
\label{supp:protocol}
\begin{itemize}
  \item DrawBench order and seeds (seed = prompt index); reference = uncached 50-step output of the same code path.
  \item Budget mapping of the official \sea gate on \flux: $\delta=0.2/0.3/0.37/0.5/0.6/0.8/1.0\to
  27/20.8/18.1/15/13/11/9$ \nfe. \seacache's Table~1 labels its $\approx$50\% row $\delta=0.3$ but its TFLOPs imply
  18.45 full steps, which matches our $\delta\approx0.37$; its $\delta=0.6$ row (12.99 steps) matches ours (13.00).
  We therefore match budgets by measured \nfe.
  \item SSIM variant (scikit-image, $7{\times}7$ uniform window) vs.\ the Gaussian SSIM of TeaCache's evaluation code;
  LPIPS (AlexNet, inputs in $[-1,1]$); PSNR on 8-bit RGB.
  \item Video: 944 unique VBench prompts (two duplicates in VBench's list of 946); frame counts (HunyuanVideo 65, Wan2.1
  81, the latter matching \seacache's reported TFLOPs); metrics on decoded MP4 frames with TeaCache's metric code;
  \TODO{seeds}.
  \item Bootstrap: paired over prompts, 10{,}000 resamples, percentile intervals, identical resampling indices for all
  pairs.
\end{itemize}

\section{Additional results}
\label{supp:results}
\begin{itemize}
  \item Full budget sweep with confidence intervals for all fillers, including linear extrapolation
  (table view of the main-paper trade-off figure plus a linear-extrapolation column).
  % source: tables/tab_flux_sweep.tex (generated); add the respect_<d>_w0 runs as 'linear extrapolation'.
  \item Open-loop forecast error inside gaps on the development set (six reference trajectories): relative velocity
  error per skipped step for the residual, clean-image, noise and velocity targets. \TODO{figure.}
  \item Correction size $\|\lat'_b-\lat_b\|/\|\lat_b\|$ versus gap length and position, all budgets. \TODO{figure.}
  \item Development-set (COCO-64) results for all configurations, with the halves used for development and
  confirmation. \TODO{development-set table.}
  % source: research_v2/ablations/ABLATIONS_devset.md
  \item Full ablation table, including the continuation variants.
  % source: tables/tab_ablation_full.tex (generated by research_v2/ablations/make_tables.py)
  \item Video: per-dimension VBench scores and per-budget tables. \TODO{after the runs.}
  \item TFLOPs and latency for all methods. \TODO{after the controlled measurement.}
\end{itemize}

\section{Negative and neutral results}
\label{supp:negative}
\begin{itemize}
  \item Forecast target: forecasting the clean-image, noise or velocity prediction reconstructed from the current
  latent, instead of the block-stack residual, is not better on the development set (all differences within about
  $\pm$0.13\,dB, intervals include zero).
  \item Allocation: an innovation-driven controller that replaces the \sea gate. \TODO{Owner decision: the interim
  development-set data show a large schedule (allocation) effect that a fixed schedule reproduces, and the predictive
  variant fails; this is not evidence that error control is worse than the \sea gate. Decide whether to report it.}
\end{itemize}

\section{More qualitative results}
\label{supp:qualitative}
\begin{itemize}
  \item Qualitative results at 13 and 20.8 \nfe with the selection rule of the main paper and more quantiles.
  % figures/fig_qualitative_supp_nfe13.pdf, figures/fig_qualitative_supp_nfe21.pdf
  \item Video frames. \TODO{after the runs.}
\end{itemize}
